\subsection{5.3 精化表示层}
\noindent 界面按角色的主要决策组织信息，使用浅雪蓝、冰湖青与暖金区分任务优先级。全角色可退出登录；所有非顾客角色有本人上下班打卡。颜色只作辅助，状态同时有文字、数量和具名责任人。移动端保持购物车、退出和核心任务入口可见。
\begin{BSRTable}
\rowcolor{TableHead}\color{white}\bfseries 角色与工作区 & \color{white}\bfseries 页面首要信息及可执行动作 \\
顾客 & 先选 Europe/Oslo 起租日 09:00 后至归还日 16:00 前的日期段，再显示整个时段可订数量；按类别、款式、尺码、滑雪水平筛选，加入购物车。订单展示 NOK 租金、每件参考价值 8\% 押金、超时规则和应还时刻；阅读损坏条款并勾选后才可结算。进度显示每一步责任岗位。\\
租赁专员 & 新单和未解决单以小红色感叹号及数量提醒，准备清单逐件绑定、复核与核码交付；准备进度同步给顾客。归还时逐件登记时间、当前门店、附件和异常，并查看具名下一处理人。\\
门店经理 & 首页突出待复核、超期初评、未解决订单和本店每日排班；仅查看本店员工名册、班次与逐件库存，另可查跨店库存聚合数量。安排本店员工班次、分派工单、审批损坏费用和安全验收；外店班次及已有打卡班次的修改由服务端拒绝。\\
维修技师 & 展示本人基本资料、资格状态、获派工单、优先级与时限；分步上传受损部位图片、初评等级、处理记录和安全检查。受损初评须在归还后 24 小时内反馈；技师不能自行批准恢复可租。\\
经营负责人 / 权限管理员 & 经营页展示订单、库存、营收和异常汇总；库存按门店/状态计数，不展示逐件编号、成本或检查备注。权限页管理账号、角色和安全配置审计。两者均可本人打卡，但不自动取得门店经理的日常写权限。\\
\end{BSRTable}
\subsubsection{5.3.1 交互与同步}
\noindent 购物车数量不能超过选定时段可订量，报价只读且不锁库存；提交时服务端原子再核。跨角色更新后，授权页面收到不含敏感内容的修订事件，再重新请求业务数据。编辑中的表单保留用户输入，版本冲突提示重读。新单标记的已读与业务未解决状态分开；后者只有完成对应处理才消失。
\noindent 日期、金额、责任人和状态不得依赖图标猜测。支付、退款、签署接入第三方之前均标明为本地演示；退款区分“已发起”与“已到账”。若受损费用超过押金，显示补差金额和专员联系状态。新顾客首笔有效租赁享 2 小时超时宽限，取消单不消耗资格。
\subsubsection{5.3.2 接口与可用性边界}
\noindent 报价、预约、准备、交付、归还、鉴定、费用审核和押金结算接口均校验当前身份、门店范围、对象版本及幂等键。顾客只看本人订单；技师只更新获派工单；经理审核与技师提交分权。输入有标签与键盘焦点，390 像素手机宽度页面不整体横向溢出；数据表可在局部滚动。错误指出库存不足、已过期、条款未同意、重复打卡或版本冲突的具体原因。
